
% ------------------------------------------------------------------
% Intro
% DSHARP spherical and glitterin irregular
% Same general assumptions
% Used to constrain a_max
% re-state scale factor and chi^2
% Mode agreement/ difference s



% ------------------------------------------------------------------
% \subsection{Overall}
% Can the models reproduce the observed polarization fractions?
% Which wavelengths are reproduced well/poorly?
% Is there an overall model that performs best?
% Does excluding Band 5 change the conclusion?



% The dust scattering models provide a method to constrain the grain properties that best replicate the observed polarization fractions of IRS 63

% For spherical grains, the model 


% ------------------------------------------------------------------



% ------------------------------------------------------------------
% \subsection{Effect of Filling Factor}
% Tested many
% 0.5 is best (porous)
% changing f changes scattering properties and a_max
% is not certain


The scattering models for spherical grains at various porosities showed that compact ($f = 1$) and very porous grains ($f = 0.01$) provide significantly poorer fits to the observations compared to the intermediate filling factors. The goodness-of-fit improves as the filling factor decreases from $f = 1$ toward $f = 0.5$, reaches a minimum at $f = 0.5$, and then worsens towards $f = 0.01$. This preference for a porous grain indicates that the scattering properties required to reproduce the observations are best represented by porous grains, rather than a fully a compact or very porous grain. This trend was seen at all fitting cases, suggesting the preferred porosity does not depend on the wavelengths included in the fit. 




% \textbf{Effect on \Pw :}From Figures \ref{fig: Spherical_BEST}, \ref{fig: appendix_all_band}, and \ref{fig: appendix_just_47} we can visually see how changing the porosity ($f$) changes the polarization curves. I will focus on Figure \ref{fig: Spherical_BEST} since that is the best fit. 
% \note{Make \Pw vs a for different porosities, without scale factor. }

% \textbf{Effect on maximum grain size: }From Table \ref{tab: dust_model_fits_spherical}, the maximum best-fit dust grain size increases as the porosity decreases, meaning that when grains are compact, the model best matches the observations at smaller grain sizes, whereas more porous grains need a larger grain size to match the observations. \add{some plot to talk about this more?} \add{model sucess vs grain size}. This shows that the porosity and success of the fit do have some sort of correlation, as the results are not random. 
% ------------------------------------------------------------------




% ------------------------------------------------------------------
% \subsection{Grain Size}
% ISM 0.1
% Dense core 1
% both models larger grains 
% compare with other disks 

The maximum grain size for also varies with filling factor. As the filling factor decreases, the best-fit maximum grain size increases, ranging from $162$ - 535 \micron for the spherical models. This trend is likely due to the lower amount of material in a porous grain. As the filling factor decreases, a larger radius may be required to produce the same scattering properties of compact grains. Irregular grain models, where porosity is not varies, therefore, has a much smaller range of best-fit grain sizes across the fitting cases, ranging from 209-212\micronNS. The best-fit maximum dust grain size is \sbest for spherical grains with $f=0.5$, and \ibest for irregular grains.

 
Figure \ref{fig: Kataoka2015Figure3Recreation} showed that the \Pweff curves for compact grains do not peak at the characteristic grain size, but rather at a slightly larger grain size. Despite this slight difference due to varying filling factor, it can still be used to estimate the grain size the observations are most sensitive to. The characteristic grain sizes that correspond to the observed wavelength range from $\sim140-340$\micron (Table \ref{tab: amax_peak_values}). The best-fit grain sizes and many of the other fitted grain sizes, fall within this range of grain sizes expected to be detected at the observed wavelengths. 

% 156.67, 115. 67, 101.58, 69 


% Al all shape, porosity, and filling combinations the maximum fit grain size ranges from 162-535 \micronNS.  The characteristic grain sizes for the observed wavelengths range from $\sim140-340$ \micronNS, indicating that the best-fit grain size, and majority of the other fitted grain sizes, are consistent within the grain sizes we expected to be able to probe. 






% ------------------------------------------------------------------













% ------------------------------------------------------------------
\subsection{Dust Grain Size Comparison}


\SCC{ There are other similar studies for disks from DSHARP.  The grain sizes from thermal dust emission tend to be larger than from scattering}

\SCC{\textbf{thermal emission study of irs 63:} You can check to see if anyone has done a grain size model for IRS 63 using thermal dust emission. So you want to discuss the grain size in IRS 63 not just from scattering and what we get from from thermal emission as well.  I can't recall if anyone modeled the thermal emission (you can see what eDisk has done) of IRS 63.  If not, you can put it in the context of what has been found for scattering for other disks and what has been found for thermal emission for other disks.  IRS 63 is younger than HL Tau, but how does its grain size distribution compare?}


----------------------------------------------------------------------------------------------------------------------------------------------------------


\cite{Kataoka_2016} Contrained the dust grain size using polarization observations for HL Tau, and found a maximum grain size of $\sim$ 150 \micronsNS. This work was extended by \cite{Yang2016} to include the polarization induced by disc inclination with respet to line of sight, and they found a value of 72 \micronNS. 

\cite{Carrasco} for HL Tau fits millimetree SEC, at the centre paricles are around a few millimetres, and maximum grain size decreases with radius (1 mm at 100 au). thermal emission?

\cite{Mori} HL Tay, 3.1 mm polarization, radiative trasfer similations, maximum grain size of $\sim$ 130 

\cite{Ueda} HL Tau $\sim$ 160




\cite{Ohashi} studies IRS 48 protoplanetary disk. This disk is also in Ophiuchus. They found $\sim$ 5 cm dust grains at the centre of the dust trap and $\sim$ 100 \micron outside them or (2) $\sim$100 \micron throughout the ring. This was using polarization.



\cite{Bacciotti_2018} uses polarization data for twwo Class II disks CW Tay and DG Tau. 100-150 \micron for CW Tay and 50-70 \micron for DG Tau. 

\cite{Hu} HH 212 found $\sim$ 130 \micron. VLA observations, model observation using radiative transfer calculation. Same disk, \cite{Lee_2021} 75 micron Band 7, 50-75 \micronNS. 


\cite{Mori} HL Tay, 3.1 mm polarization, radiative trasfer similations, maximum grain size of $\sim$ 130 
% ------------------------------------------------------------------




% ------------------------------------------------------------------
\subsection{Spherical vs Irregular Grains}

The spherical and irregular grain models allow us to determine the role of grain shape in polarized scattering. The lowest goodness-of-fit parameter $\chi_\mathcal{P}^2$ for spherical and irregular grain models is 1.37 and 4.54, respectively. A lower $\chi_\mathcal{P}^2$ value indicates that the best-fit spherical model matches the observational polarization fractions compared to irregular grains. However, this does not provide strong enough evidence to claim that IRS has spherical grains over irregular grains. For example, the compact spherical grain model fit has a much higher $\chi_\mathcal{P}^2$ value of $28.80$ compared to 4.54 for the irregular grain model. This suggests that the grain shape and porosity both influence the polarization from scattering, and we cannot constrain the shape of the grain from these observations. 




% ------------------------------------------------------------------



 
% ------------------------------------------------------------------
\section{Implications for Dust Evolution}

The best-fit maximum grain sizes of \sbest and \ibest, for spherical and irregular grains respectively, are significantly larger than grain sizes expected in the ISM and the dense core of the molecular cloud, \ism and \dccNS, respectively. This shows evidence that dust grain growth is occurring in IRS 63. 

While this result is relevant and important to the disk itself, it has broader implications for planet formation as a field. This suggests the earliest steps of planet formation are occurring at this early evolutionary stage of a disk (Class I), as predicted by \cite{Dom2020}. This does not suggest that planets have formed by this stage, but rather that planet formation might begin earlier than some planet formation models might consider. Studies of the earlier stages of protostellar disks have a great impact on the field of planet formation as they probe these earlier stages.




% ------------------------------------------------------------------